\begin{definition}[Physical incidence coordinates of an open region]%
    \label{def:peps_region_half_edge_label_equiv}
    \lean{TNLean.PEPS.regionHalfEdgeLabelEquiv}
    \leanok
    For a region $R$ in a finite simple graph and an arbitrary alphabet $X$,
    regroup the physical incidence labels as one label on each crossing bond,
    one tail label on each internal bond, and one head label on each internal
    bond. This is a genuine coordinate bijection. The construction uses the
    native ordered endpoints and retains every incidence at the boundary.
    It is the incidence part of the regional coordinates in
    \cite[Section~6.2]{Schuch2010PEPS} and of the bond-endpoint argument in
    \cite[Section~7]{Schuch2010PEPS}.
\end{definition}

\begin{lemma}[Actual endpoint coordinates in the region bijection]%
    \label{lem:peps_region_half_edge_label_endpoints}
    \lean{TNLean.PEPS.regionHalfEdgeLabelEquiv_apply_internal_tail,
        TNLean.PEPS.regionHalfEdgeLabelEquiv_apply_internal_head,
        TNLean.PEPS.regionHalfEdgeLabelEquiv_apply_boundary_tail,
        TNLean.PEPS.regionHalfEdgeLabelEquiv_apply_boundary_head}
    \leanok
    \uses{def:peps_region_half_edge_label_equiv}
    The two internal labels are precisely the physical labels at the ordered
    tail and head. On a crossing bond, the single retained physical label is
    the tail label when the tail lies in $R$, and the head label when the
    head lies in $R$.
\end{lemma}

\begin{proof}\leanok
    Each internal edge has exactly two incidences in $R$, whereas each crossing
    edge has exactly one. Read these incidences in the coordinate bijection.
\end{proof}
